\chapter{Portfolio Construction}\label{ch:iv2:portfolio-construction}

\begin{ivtrains}
Building a portfolio by hand, from the two-asset case to constraints, risk parity, factor models, the fundamental law, shrinkage and turnover. Returns are annual and in excess of cash unless a stem says otherwise, and every number is invented.
\end{ivtrains}

\begin{question}[$\star$ \iqroles{QR, QT, Analyst} \iqfirm{any} \iqlevel{junior}]\label{iq:iv2:portfolio-construction:1}
Two stocks have volatilities of 20\% and 30\%, and a correlation of 0.5. What weights give the minimum-variance portfolio? What is its volatility?
\end{question}

\begin{question}[$\star$ \iqroles{QR, QT, Analyst} \iqfirm{any} \iqlevel{junior}]\label{iq:iv2:portfolio-construction:2}
You have two assets with the same volatility of 20\%, and a correlation of 0.2. What is the minimum-variance mix, and its volatility?
\end{question}

\begin{question}[$\star$ \iqroles{QR, Risk, Analyst} \iqfirm{any} \iqlevel{junior}]\label{iq:iv2:portfolio-construction:3}
Two assets have volatilities of 10\% and 20\%. What weights make their risk contributions equal? Does the answer depend on the correlation?
\end{question}

\begin{question}[$\star$ \iqroles{QR, QT} \iqfirm{hedge fund} \iqlevel{junior}]\label{iq:iv2:portfolio-construction:4}
A signal has an information coefficient of 0.05. Using the fundamental law, what is the information ratio with 100 independent bets a year? With 252?
\end{question}

\begin{question}[$\star$ \iqroles{QR, QT, Trader} \iqfirm{hedge fund, market maker} \iqlevel{junior}]\label{iq:iv2:portfolio-construction:5}
A portfolio moves from weights (0.5, 0.3, 0.2) to (0.4, 0.4, 0.2). What is the one-way turnover, and what do you pay at 10 basis points per dollar traded? If you do this every month, what does it cost per year?
\end{question}

\begin{question}[$\star$ \iqroles{QT, Trader, Risk} \iqfirm{hedge fund} \iqlevel{junior}]\label{iq:iv2:portfolio-construction:6}
You are long \$10 million of a stock with a beta of 1.2. You hedge with a stock with a beta of 0.8. How much do you short for beta neutrality? What is your net dollar position?
\end{question}

\begin{question}[$\star\star$ \iqroles{QR, QT} \iqfirm{hedge fund, bank} \iqlevel{junior, experienced}]\label{iq:iv2:portfolio-construction:7}
Two assets have expected excess returns of 8\% and 10\%, variances of 0.04 and 0.09, and a covariance of 0.03. What are the tangency weights and its Sharpe ratio? For a mean-variance investor with a risk aversion of 4, what do they hold, and what are its return and volatility?
\end{question}

\begin{question}[$\star\star$ \iqroles{QR, QT, Risk} \iqfirm{any} \iqlevel{junior, experienced}]\label{iq:iv2:portfolio-construction:8}
Two stocks have volatilities of 20\% and 30\% and a correlation of 0.8. What is the minimum-variance portfolio? Why is one weight negative? Give the exact condition.
\end{question}

\begin{question}[$\star\star$ \iqroles{QR, QT} \iqfirm{hedge fund} \iqlevel{junior, experienced}]\label{iq:iv2:portfolio-construction:9}
Asset A has an excess return of 10\% and a volatility of 20\%. Asset B has 5\% and 30\%. Their covariance is 0.05. Find the Sharpe ratio of the unconstrained tangency portfolio, and give its weights as $\Sigma^{-1}\mu$, the holdings at a risk aversion of 1. What is the best long-only Sharpe, and how do you know?
\end{question}

\begin{question}[$\star\star$ \iqroles{QR, Risk, Analyst} \iqfirm{any} \iqlevel{junior}]\label{iq:iv2:portfolio-construction:10}
Stocks each have 30\% volatility and a common pairwise correlation of 0.3. What is the volatility of an equally weighted portfolio of 10, 30 and 100 of them? What is the limit? How many stocks do you need to be within 1 point of that limit?
\end{question}

\begin{question}[$\star\star$ \iqroles{QR, QT} \iqfirm{hedge fund} \iqlevel{junior, experienced}]\label{iq:iv2:portfolio-construction:11}
You rank 500 stocks every year with an information coefficient of 0.02, and the forecast errors have a pairwise correlation of 0.01. What is the effective breadth and the information ratio? Compare it with what the fundamental law gives if you treat them as independent.
\end{question}

\begin{question}[$\star\star$ \iqroles{QR, QT} \iqfirm{hedge fund} \iqlevel{junior, experienced}]\label{iq:iv2:portfolio-construction:12}
A signal has an information coefficient of 0.05 and 500 independent bets. The transfer coefficient of the constrained long-only portfolio is 0.3. What is the expected information ratio? What does the coefficient measure, and how would you improve it?
\end{question}

\begin{question}[$\star\star$ \iqroles{QR, QT, Analyst} \iqfirm{hedge fund, bank} \iqlevel{junior, experienced}]\label{iq:iv2:portfolio-construction:13}
Two assets have sample mean excess returns of 15\% and $-5\%$, from 20 years of data, with volatility 20\% each. What is the standard error of each mean? Shrink both halfway to a common 5\%. What happens? Why does mean-variance optimisation make estimation error worse?
\end{question}

\begin{question}[$\star\star$ \iqroles{QR, Risk} \iqfirm{hedge fund, bank} \iqlevel{junior, experienced}]\label{iq:iv2:portfolio-construction:14}
A sample covariance has eigenvalues 5, 1 and 0.01. What is its condition number? Shrink it 20\% to a multiple of the identity with the same trace. What are the new eigenvalues and condition number? Why does this help a minimum-variance optimiser?
\end{question}

\begin{question}[$\star\star$ \iqroles{QR, QT} \iqfirm{hedge fund} \iqlevel{junior, experienced}]\label{iq:iv2:portfolio-construction:15}
You compute a correlation matrix for 100 stocks from 1,000 days of pure noise. Where do the eigenvalues lie, by Marchenko-Pastur? What happens with 500 stocks? How do you use that?
\end{question}

\begin{question}[$\star\star$ \iqroles{QR, Risk, QD} \iqfirm{hedge fund, bank} \iqlevel{junior, experienced}]\label{iq:iv2:portfolio-construction:16}
For 500 stocks, how many parameters does a full covariance matrix have? How many does a model have with 10 factors, a full factor covariance and individual idiosyncratic variances? What assumption buys you that, and when does it fail?
\end{question}

\begin{question}[$\star\star$ \iqroles{QR, QT, Analyst} \iqfirm{any} \iqlevel{junior, experienced}]\label{iq:iv2:portfolio-construction:17}
Two independent assets each have a zero expected arithmetic return and 20\% volatility, lognormal. A portfolio is rebalanced yearly to 50/50. What is its expected log growth, against each asset's? How does buying and holding compare over 30 years? What if the correlation is 1?
\end{question}

\begin{question}[$\star\star\star$ \iqroles{QR, QT} \iqfirm{hedge fund} \iqlevel{experienced}]\label{iq:iv2:portfolio-construction:18}
Ten assets each have a true mean of 0.5\% a month, volatility of 5\%, and pairwise correlation of 0.3. You estimate the tangency portfolio from 60 months of data. What is the true Sharpe of equal weights? Roughly what in-sample and out-of-sample annualised Sharpe do you get from the estimated weights? What do you take from it?
\end{question}

\begin{question}[$\star\star\star$ \iqroles{QR, QT, Analyst} \iqfirm{hedge fund, bank} \iqlevel{experienced}]\label{iq:iv2:portfolio-construction:19}
Market weights are 0.6 and 0.4, the variances are 0.04 and 0.09, the covariance is 0.03 and risk aversion is 2.5. What are the implied excess returns? You have the view that asset 1 returns 10\%, with $\tau = 0.05$ and view uncertainty equal to $\tau$ times its variance. Give the Black-Litterman posterior returns and weights. What happens as the uncertainty goes to zero?
\end{question}

\begin{question}[$\star\star\star$ \iqroles{QR, Risk} \iqfirm{hedge fund, bank} \iqlevel{junior, experienced}]\label{iq:iv2:portfolio-construction:20}
Three assets have variances 0.04, 0.01 and 0.09. The covariance is 0.01 between the first two, 0.004 between the last two, and 0 between the first and the third. Compute the risk contribution shares for equal weights, for inverse-volatility weights, each proportional to one over the asset's volatility, and for (0.2763, 0.5119, 0.2117). Which is the equal-risk portfolio? What does it have that inverse volatility does not?
\end{question}

\begin{question}[$\star\star\star$ \iqroles{QR, QD} \iqfirm{any} \iqlevel{experienced}]\label{iq:iv2:portfolio-construction:21}
Three assets have expected excess returns of 6\%, 8\% and 12\%, and variances 0.04, 0.09 and 0.0225. Their covariances are 0.012, 0.012 and 0.018, for pairs (1, 2), (1, 3) and (2, 3). Find the global minimum-variance portfolio's return and volatility. Find the portfolio on the frontier with a return of 12\%. How does the two-fund theorem show up here?
\end{question}

\begin{question}[$\star\star\star$ \iqroles{QR, QD, Risk} \iqfirm{any} \iqlevel{experienced}]\label{iq:iv2:portfolio-construction:22}
Three assets have variances 0.04, 0.09 and 0.0225, and covariances 0.012, 0.012 and 0.018 for pairs (1, 2), (1, 3) and (2, 3). With weights summing to 1, their unconstrained minimum-variance weights are (0.268, 0.043, 0.689). Impose a maximum weight of 40\%. What is the solution? How do you know it is optimal, by the conditions and not just by a solver?
\end{question}

\begin{question}[$\star\star\star$ \iqroles{QR, QT, Trader} \iqfirm{hedge fund, market maker} \iqlevel{junior, experienced}]\label{iq:iv2:portfolio-construction:23}
A signal says 2 basis points a day of expected return, decaying with a half-life of 5 days. Trading costs 5 basis points each way. Is the trade worth doing? What is the break-even starting signal? What changes if the position must be sold after 1 day?
\end{question}

\begin{question}[$\star\star\star$ \iqroles{QR, QT} \iqfirm{hedge fund} \iqlevel{junior, experienced}]\label{iq:iv2:portfolio-construction:24}
Two strategies have Sharpe ratios of 1.0 and 0.8, and a correlation of 0.3. What is the Sharpe of the best mix, and with what risk allocation? At what correlation does the second strategy add nothing? What happens at a negative correlation?
\end{question}

\begin{question}[$\star\star\star\star$ \iqroles{QR, QT, Risk} \iqfirm{hedge fund} \iqlevel{experienced}]\label{iq:iv2:portfolio-construction:25}
You have three stocks with expected returns of 6\%, 3\% and $-4\%$, volatilities of 25\%, 20\% and 30\%, common correlation of 0.4, and betas of 1.2, 0.8 and 1.0. Your risk aversion is 5. Find the unconstrained portfolio and its beta. Then make it beta neutral, and then dollar neutral as well. Compare return, volatility and Sharpe. Scale to a 10\% volatility. Then add a cost of 0.2\% per dollar traded starting from flat. What changes, and how do you check the answer without a solver?
\end{question}
